% Zarathustra: Learned Generative Modeling of I/O Workload Traces
%
% Submission target: USENIX FAST / ACM SOSP / IEEE MASCOTS
% Compile: pdflatex main && bibtex main && pdflatex main && pdflatex main

\documentclass[10pt,twocolumn]{article}

%% ---------- packages ----------
\usepackage[margin=1in]{geometry}
\usepackage{times}
\usepackage{microtype}
\usepackage{amsmath,amssymb}
\usepackage{graphicx}
\usepackage{booktabs}
\usepackage{xspace}
\usepackage{url}
\usepackage{hyperref}
\usepackage{xcolor}
\usepackage{listings}
\usepackage{subcaption}
\usepackage{multirow}
\usepackage{siunitx}

%% ---------- macros ----------
\newcommand{\sys}{Zarathustra\xspace}
\newcommand{\eg}{\emph{e.g.,}\xspace}
\newcommand{\ie}{\emph{i.e.,}\xspace}
\newcommand{\etal}{\emph{et al.}\xspace}
\newcommand{\todo}[1]{\textcolor{red}{\textbf{[TODO: #1]}}}
\newcommand{\mmd}{MMD$^2$\xspace}
\newcommand{\dmdgen}{DMD-GEN\xspace}

%% ---------- metadata ----------
\title{\textbf{\sys: Learned Generative Modeling of\\Storage I/O Workload Traces}}

\author{
  Darrell D.\ Long \\
  University of California, Santa Cruz \\
  \texttt{darrell@ucsc.edu}
  \and
  \todo{co-authors}
}

\date{}

\begin{document}

\maketitle

%% ============================================================
\begin{abstract}
%% ============================================================
Storage system evaluation requires realistic I/O workload traces, but
production traces are proprietary, voluminous, and tied to a single
point in time.  \sys is a generative model that learns the full joint
distribution of I/O request arrivals, object sizes, access patterns,
and read/write ratios from a corpus of real storage traces.  We build
on the LSTM-GAN (LLGAN) architecture~\cite{zhang2024llgan} and extend
it with a latent autoencoder with autoregressive supervisor training
(following TimeGAN~\cite{yoon2019timegan} and SeriesGAN~\cite{seriesgan2024}),
a Wasserstein critic with spectral normalization, locality-aware
object-ID representation, and a suite of auxiliary losses targeting
frequency content and distributional moments.  Training on hundreds
of publicly available Tencent and Alibaba block storage traces, \sys
generates synthetic traces that are statistically indistinguishable
from real production data across \todo{N} metrics, including
MMD$^2$, PRDC precision/recall, Context-FID, AutoCorr, and DMD-GEN.

\end{abstract}


%% ============================================================
\section{Introduction}
\label{sec:intro}
%% ============================================================

Storage system research depends on representative I/O workloads for
benchmarking cache policies, capacity planning, and system stress
testing.  Real production traces capture the full complexity of
workload behavior---burstiness, temporal locality, mixed request sizes,
and correlated access patterns---but are rarely available: they are
proprietary, enormous, and snapshot-in-time.

Synthetic trace generators have long been proposed as an alternative,
but classical approaches rely on hand-tuned Markov chains or
parameterized distributions that fail to capture the joint statistical
structure of real workloads~\cite{paul2022hpc,geminio2024}.  Recent
work has shown that deep generative models can learn this structure
from real data~\cite{zhang2024llgan}, but existing systems train on a
single trace and fail to generalize across workloads.

We present \sys, a generative model trained on a corpus of hundreds
of storage traces that produces statistically realistic synthetic I/O
sequences of arbitrary length.  Our contributions are:

\begin{enumerate}
  \item A latent autoencoder + adversarial training curriculum
        (Section~\ref{sec:architecture}) that improves mode coverage
        and temporal coherence over the LLGAN baseline.

  \item A locality-aware object-ID representation using signed
        log-delta encoding that captures repeat accesses, sequential
        strides, and random jumps within a bounded tanh output range
        (Section~\ref{sec:preprocessing}).

  \item A suite of auxiliary losses---frequency-domain spectral loss
        with FIDE amplitude weighting~\cite{fide2024}, moment matching,
        and feature matching~\cite{r3gan2024}---that constrain the
        generator beyond the adversarial signal alone
        (Section~\ref{sec:training}).

  \item A comprehensive evaluation framework
        (Section~\ref{sec:eval}) incorporating \mmd, PRDC
        precision/recall/density/coverage, Context-FID, autocorrelation
        score, and DMD-GEN~\cite{dmdgen2025}.

  \item An empirical study of multi-corpus generalization across
        Tencent Block 2020 and Alibaba Block 2020 trace corpora
        (Section~\ref{sec:results}).
\end{enumerate}


%% ============================================================
\section{Background and Related Work}
\label{sec:background}
%% ============================================================

\subsection{Storage Trace Generation}

Classical trace generation fits parametric models (\eg Poisson
arrivals, Pareto size distributions) to observed statistics.  These
models reproduce marginal distributions but lose the joint temporal
structure---burstiness, locality, phase transitions---that drives
cache behavior.  LLGAN~\cite{zhang2024llgan} showed that LSTM-based
GANs can reproduce block I/O statistics more faithfully than
parametric models.  \sys extends this to multi-file corpus
generalization and a substantially improved training objective.

\subsection{Time Series Generation}

TimeGAN~\cite{yoon2019timegan} introduced a latent autoencoder +
supervisor curriculum for time series generation that avoids the
instability of direct adversarial training in feature space.
SeriesGAN~\cite{seriesgan2024} extended this with a dual discriminator
and two-step supervisor, reporting a 34\% reduction in discriminative
score over TimeGAN.  AVATAR~\cite{avatar2025} replaces the recovery
network with an adversarial autoencoder posterior, simplifying the
latent space alignment.  ChronoGAN~\cite{chronogan2024} adds
slope and skewness terms to the time-series loss.  We draw from
all three architectures but adapt them to the mixed discrete/continuous
structure of I/O traces.

\subsection{GAN Objectives and Regularization}

Wasserstein GANs~\cite{gulrajani2017wgangp} stabilize training by
bounding the critic's Lipschitz constant.  R3GAN~\cite{r3gan2024}
proposes a relativistic paired loss (RpGAN) with joint R1/R2
gradient regularization and shows improved mode coverage on
StackedMNIST.  Spectral normalization~\cite{miyato2018spectralnorm}
enforces the Lipschitz constraint without second-order autograd and
is compatible with Apple MPS hardware.

\subsection{Frequency and Distributional Losses}

Diffusion-TS~\cite{diffusionts2024} demonstrates FFT-domain spectral
losses for time series generation.  FIDE~\cite{fide2024} introduces
frequency inflation weighting to prevent rare extreme values from being
washed out by a flat spectral MSE.  We incorporate both.

\subsection{Evaluation Metrics}

MMD$^2$~\cite{todo} detects distributional mismatch but misses
mode collapse.  PRDC~\cite{naeem2020reliable} adds precision, recall,
density, and coverage.  TSGBench~\cite{tsgbench2024} defines a
12-metric benchmark for time series generation; we adopt Context-FID
and autocorrelation score from it.  DMD-GEN~\cite{dmdgen2025}
quantifies temporal dynamics fidelity via Grassmannian comparison of
dynamic mode decomposition eigenvectors; it detects autocorrelation
and burst-structure failures that MMD$^2$ and PRDC miss.

\subsection{Path-Space Methods}

PCF-GAN~\cite{pcfgan2023} defines discrimination in path space via the
characteristic function of measures on the path space, providing
stronger sensitivity to sequential law than hidden-state pooling.
The High-Rank Path Development method~\cite{hrpcf2024} extends this
to avoid rank collapse.  Sig-WGAN~\cite{sigwgan2021} replaces the
LSTM critic with a path-signature-based Wasserstein distance.  These
approaches are stronger theoretical foundations for a future critic
redesign.


%% ============================================================
\section{Problem Formulation}
\label{sec:problem}
%% ============================================================

An I/O trace is a sequence of requests
$\mathbf{x} = (x_1, x_2, \ldots, x_N)$ where each event
$x_t = (\Delta\tau_t, o_t, s_t, c_t, u_t)$ contains:
\begin{itemize}
  \item $\Delta\tau_t$: inter-arrival time (seconds) to the previous request
  \item $o_t$: object (block) identifier
  \item $s_t$: object size (bytes)
  \item $c_t \in \{+1,-1\}$: operation code (read/write)
  \item $u_t$: tenant identifier
\end{itemize}

We model traces in a normalized feature space.  Let
$\mathbf{X} = \text{prep}(\mathbf{x}) \in [-1,1]^{N \times d}$
be the preprocessed sequence ($d = 5$).  The generative model
$G_\theta$ produces sequences $\hat{\mathbf{X}} \sim G_\theta(z)$
from noise $z$, trained so that $\hat{\mathbf{X}}$ is
statistically indistinguishable from real windows
$\mathbf{X}_{t:t+T}$ of length $T$.

At generation time, the inverse transform
$\text{prep}^{-1}(\hat{\mathbf{X}})$ recovers absolute-time,
raw-identifier sequences suitable for replay through a storage
simulator.


%% ============================================================
\section{Preprocessing}
\label{sec:preprocessing}
%% ============================================================

\paragraph{Delta encoding.}
Raw timestamps are converted to inter-arrival times:
$\Delta\tau_t = \tau_t - \tau_{t-1}$.  This makes the series
stationary and removes dependence on the absolute start time of
each trace file.

\paragraph{Log transform.}
$\Delta\tau_t$ and $s_t$ span many decades.  We apply
$\log(1 + x)$ before min-max normalization to $[-1,1]$,
compressing the dynamic range while preserving monotonicity.

\paragraph{Locality-aware object-ID encoding.}
Raw $o_t$ values cannot be regressed directly: they span up to
$10^{10}$ with no metric structure meaningful to the model.
We instead encode the signed inter-request delta:
\[
  \delta_t = o_t - o_{t-1}, \qquad
  \tilde{o}_t = \text{sign}(\delta_t) \cdot \log(1 + |\delta_t|).
\]
This representation encodes locality structure natively:
$\tilde{o}_t \approx 0$ is a cache hit or repeat access;
small $|\tilde{o}_t|$ is sequential or strided access;
large $|\tilde{o}_t|$ is a random jump.  At generation time,
$o_t = o_0 + \text{cumsum}(\text{undo-signed-log}(\hat{\tilde{o}}))$
reconstructs absolute identifiers.

\paragraph{Opcode binarization.}
Read $\to +1$, write $\to -1$, aligned with the tanh output range.

\paragraph{Multi-file normalization.}
The preprocessor is fit once on a seed sample of files and held
fixed across all training epochs and at generation time.


%% ============================================================
\section{Architecture}
\label{sec:architecture}
%% ============================================================

\sys uses a latent autoencoder GAN with four components:

\begin{description}
  \item[Encoder $E$] maps real feature windows
    $\mathbf{X} \in \mathbb{R}^{T \times d}$ to latent sequences
    $\mathbf{H} = E(\mathbf{X}) \in \mathbb{R}^{T \times m}$.
    Single-layer LSTM, hidden size $m$.

  \item[Recovery $R$] decodes latent sequences back to feature space:
    $\hat{\mathbf{X}} = R(\mathbf{H}) \in \mathbb{R}^{T \times d}$.
    Single-layer LSTM + linear projection.

  \item[Supervisor $S$] models the autoregressive dynamics of the
    latent space: $S(\mathbf{H})_{t} \approx \mathbf{H}_{t+k}$
    for step size $k \in \{1, 2\}$ (2-step following SeriesGAN).
    Single-layer LSTM.

  \item[Generator $G$] maps noise to latent sequences.
    A global noise vector $z_g \in \mathbb{R}^{n}$ is projected to the
    LSTM initial hidden state (workload identity); per-step local noise
    $z_l \in \mathbb{R}^{T \times n}$ provides per-step innovation.
    Output: $\mathbf{H}_{\text{fake}} = G(z_g, z_l) \in [0,1]^{T \times m}$
    (sigmoid-bounded to match the latent AE range).

  \item[Critic $C$] scores latent sequences as real or generated.
    Single-layer LSTM with learned attention pooling over hidden states
    (replacing mean pooling to focus on burst and regime-change
    timesteps).  Output linear layer uses spectral normalization to
    enforce the Wasserstein Lipschitz constraint.
\end{description}

\subsection{Split Noise Latent}

The split $z_g / z_l$ design gives the generator two degrees of freedom:
$z_g$ controls the global workload character (read-heavy vs.\ write-heavy,
bursty vs.\ smooth) while $z_l$ introduces per-step stochasticity.
At generation time, $z_g$ is sampled once per stream and held fixed
across all windows, while a new $z_l$ is drawn each window.


%% ============================================================
\section{Training}
\label{sec:training}
%% ============================================================

Training follows a four-phase curriculum:

\paragraph{Phase 1: Autoencoder pretraining.}
Optimize $E$ and $R$ with reconstruction loss
$\mathcal{L}_{\text{AE}} = \|R(E(\mathbf{X})) - \mathbf{X}\|_2^2$
for $N_{\text{AE}}$ epochs.

\paragraph{Phase 2: Supervisor pretraining.}
With $E$ fixed, optimize $S$ with
\[
  \mathcal{L}_S = \|S(E(\mathbf{X}))_{:,:-k} - E(\mathbf{X})_{:,k:}\|_2^2
\]
for $k = 2$ (two-step supervisor~\cite{seriesgan2024}).

\paragraph{Phase 2.5: Generator warm-up.}
Train $G$ with supervisor consistency loss only, before introducing the
critic.  This prevents the critic from saturating before $G$ has
learned the latent dynamics~\cite{yoon2019timegan}.

\paragraph{Phase 3: Joint GAN training.}
Alternate between:

\begin{enumerate}
  \item \textbf{Critic step} ($n_{\text{critic}}$ times): minimize
    Wasserstein loss $C(\mathbf{H}_{\text{fake}}) - C(\mathbf{H}_{\text{real}})$
    with spectral norm Lipschitz constraint.

  \item \textbf{Generator step}: minimize
    \[
      \mathcal{L}_G = -C(G(z_g,z_l))
      + \lambda_S \mathcal{L}_S^G
      + \lambda_{\text{FM}} \mathcal{L}_{\text{FM}}
      + \lambda_V \mathcal{L}_V
      + \lambda_{\text{FFT}} \mathcal{L}_{\text{FFT}}
    \]
    where $\mathcal{L}_S^G$ is the supervisor consistency loss on fake
    latents, $\mathcal{L}_{\text{FM}}$ is feature matching against the
    critic's intermediate representation~\cite{salimans2016improved},
    $\mathcal{L}_V$ is moment matching (mean + std per feature), and
    $\mathcal{L}_{\text{FFT}}$ is the FIDE-weighted spectral loss.

  \item \textbf{Encoder--Recovery step}: fine-tune $E$ and $R$ with
    reconstruction loss to keep the autoencoder aligned with the
    evolving generator distribution.
\end{enumerate}

\subsection{FIDE Spectral Loss}

The frequency-domain loss uses amplitude-weighted MSE following
FIDE~\cite{fide2024}:
\[
  \mathcal{L}_{\text{FFT}} =
    \mathbb{E}\bigl[w_f \cdot |\hat{X}_f - X_f|^2\bigr], \quad
  w_f = 1 + \alpha \frac{|X_f|}{\bar{|X|}}
\]
where $X_f$ and $\hat{X}_f$ are the real FFT magnitudes of real and
fake sequences along the time axis.  The weighting $w_f$ amplifies
loss at high-amplitude frequency bins, preventing the model from
averaging away rare burst events.


%% ============================================================
\section{Multi-File Streaming Training}
\label{sec:multifile}
%% ============================================================

Each epoch samples a random subset of $k$ files from the corpus.
Within each file, non-overlapping windows of length $T$ are extracted.
A \texttt{ConcatDataset} of per-file \texttt{TraceDataset}s ensures
no window ever crosses a file boundary---cross-file sequences would
introduce false temporal dependencies between unrelated workload
phases.

A fixed held-out set of files is removed from the training pool
before the first epoch and used exclusively for validation.  This
ensures that validation MMD$^2$ reflects generalization to unseen
files rather than in-distribution memorization.


%% ============================================================
\section{Generation}
\label{sec:generation}
%% ============================================================

At inference time, \sys generates synthetic traces of arbitrary length
by maintaining the LSTM hidden state across windows.  For a stream of
$N$ records:

\begin{enumerate}
  \item Sample $z_g$ once (workload identity).
  \item For each window: sample $z_l$, run $G(z_g, z_l, h_{\text{prev}})$,
        decode via $R$, carry $(h_n, c_n)$ to the next window.
  \item Apply inverse preprocessing per stream independently (cumulative
        sums for timestamps and object IDs must not cross stream boundaries).
\end{enumerate}

Multiple independent streams can be generated in parallel by
vectorizing over $z_g$ and $h$.


%% ============================================================
\section{Evaluation Metrics}
\label{sec:eval}
%% ============================================================

We evaluate synthetic trace quality on six complementary metrics:

\begin{description}
  \item[\mmd] Maximum Mean Discrepancy~\cite{gretton2012kernel} between
    real and generated window distributions (RBF kernel, flattened windows).
    Lower is better; zero indicates indistinguishable distributions.

  \item[PRDC] Precision, recall, density, and coverage~\cite{naeem2020reliable}
    measured in a shared feature space.  Recall captures mode coverage;
    precision captures fidelity.

  \item[Context-FID] Fréchet distance in the trained encoder's latent
    space.  Sensitive to temporal structure within windows.

  \item[AutoCorr] Average absolute mismatch between real and generated
    autocorrelation functions at lags 1--5, averaged over all features.

  \item[\dmdgen] Grassmannian distance between real and generated
    dynamic mode decomposition eigenvector subspaces~\cite{dmdgen2025}.
    Detects wrong burst structure and autocorrelation patterns that
    MMD$^2$ and PRDC miss.
\end{description}


%% ============================================================
\section{Experimental Results}
\label{sec:results}
%% ============================================================

\subsection{Dataset}

We train on the libCacheSim cache-dataset corpus~\cite{yang2021libcachesim,yang2023fifo},
comprising:
\begin{itemize}
  \item \textbf{Tencent Block 2020}: 4{,}995 volumes, 9 days, 152\,GB
    (oracleGeneral format), 512\,B--32\,KB requests, 83--94\% reads.
  \item \textbf{Alibaba Block 2020}: 1{,}001 volumes, 31 days, 93\,GB,
    4\,KB-aligned, mixed read ratios.
\end{itemize}

\subsection{Hyperparameters}

\begin{table}[h]
\centering
\small
\begin{tabular}{ll}
\toprule
Parameter & Value \\
\midrule
Window length $T$ & 12 \\
Hidden size $m$ & \todo{} \\
Noise dim $n$ & 32 \\
Latent dim & 24 \\
Batch size & 64 \\
Files per epoch & 8 \\
Records per file & 15{,}000 \\
Critic steps $n_{\text{critic}}$ & 5 \\
LR (G) & $10^{-4}$ \\
LR (C) & $10^{-4}$ \\
LR (E, R) & $5 \times 10^{-4}$ \\
$\lambda_S$ (supervisor) & 10.0 \\
$\lambda_{\text{FM}}$ (feature matching) & 1.0 \\
$\lambda_V$ (moment) & 0.1 \\
$\lambda_{\text{FFT}}$ (spectral) & 0.05 \\
FIDE $\alpha$ & 1.0 \\
Supervisor steps $k$ & 2 \\
\bottomrule
\end{tabular}
\caption{Training hyperparameters.}
\label{tab:hyperparams}
\end{table}

\subsection{Ablation Study}

\todo{Table comparing v1--v11 checkpoints across all six metrics.
Key findings expected: 2-step supervisor improves DMD-GEN;
FIDE improves FFT loss; feature matching fixes mode collapse.}

\subsection{Comparison with Baselines}

\todo{Compare against: (1) LLGAN baseline~\cite{zhang2024llgan};
(2) TimeGAN~\cite{yoon2019timegan}; (3) SeriesGAN~\cite{seriesgan2024};
(4) classical Markov-chain generator.}

\subsection{Generation Fidelity}

\todo{Figures: marginal distribution comparisons per feature;
autocorrelation plots; DMD eigenvalue spectrum comparison;
PRDC recall vs.\ training epoch.}

\subsection{Generation Length Robustness}

\todo{Rolling MMD² and AutoCorr vs.\ chunk index for streams
of 10K, 100K, 1M records.  Characterize ``max reliable stream length''.}


%% ============================================================
\section{Discussion}
\label{sec:discussion}
%% ============================================================

\paragraph{Limitations.}
The current model operates on windows of length $T = 12$, which
captures micro-burst patterns but not multi-hour workload phases or
diurnal cycles.  Hierarchical generation~\cite{stagediff2025} or
increased context via patch embedding~\cite{ttsgan2022} are natural
next steps.

The global noise vector $z_g$ currently provides stochastic workload
identity but no semantic structure.  Per-file workload
descriptors (read ratio, mean inter-arrival time, size entropy, repeat
rate) could be computed at load time and injected via FiLM
conditioning~\cite{perez2018film} to enable controllable generation
across workload profiles~\cite{constrainedts2023}.

Context-FID uses the model's own encoder as the feature extractor,
which may be co-adapted with the generator and partially self-serving
as a selection metric.  A frozen encoder trained on real data only
would provide a more objective criterion.

\paragraph{Future Work.}
A path-space critic~\cite{pcfgan2023,hrpcf2024} is a promising
replacement for the LSTM attention pooling critic, with stronger
theoretical sensitivity to sequential law and locality dynamics.
Retrieval-augmented conditioning~\cite{retrieval_diffusion2024} could
improve fidelity on statistically rare but operationally important
workload phases such as heavy-burst or hotspot episodes.
Discrete tokenization of categorical columns~\cite{sdformer2024}
would separate the modeling paths for opcode and tenant from the
continuous time and size channels.


%% ============================================================
\section{Conclusion}
\label{sec:conclusion}
%% ============================================================

\sys learns the joint distribution of storage I/O workload traces
from a corpus of real production data and generates statistically
realistic synthetic traces of arbitrary length.  Key contributions
include a latent autoencoder + supervisor training curriculum, a
locality-aware object-ID representation, FIDE-weighted spectral loss,
and a comprehensive multi-metric evaluation framework.  On the
Tencent Block 2020 corpus, \sys achieves \todo{state results}.

Realistic synthetic traces enable storage system evaluation without
requiring access to proprietary production workloads, supporting
reproducible benchmarking, stress testing, and capacity planning
across a wide range of workload regimes.


%% ============================================================
\bibliographystyle{plain}
\bibliography{refs}
%% ============================================================

\end{document}
